\documentclass{beamer}
\usetheme{Madrid}

% UTF-8 Schriftkodierung
\usepackage[utf8]{inputenc}
\usepackage[T1]{fontenc}
\usepackage[ngerman]{babel}

\usepackage{multicol}
\usepackage{tabulary}
\usepackage{booktabs}

\usepackage{soul,soulutf8}
\definecolor{pink}{rgb}{1,0.5,0.5}
\definecolor{LightSkyBlue}{rgb}{0.5,0.5,1}
\definecolor{MidnightBlue}{rgb}{0,0,0.5}
\definecolor{darkBlue}{rgb}{0,0,0.4}
\definecolor{paleBlue}{rgb}{0.8,0.8,1}
\definecolor{darkGreen}{rgb}{0,0.4,0}
\definecolor{Terminus}{rgb}{0,0.4,0}
\definecolor{Definition}{rgb}{1,0.5,0.5}
\definecolor{shade1}{rgb}{1,0.5,0.5}
\definecolor{shade2}{rgb}{1,0.5,0.5}
\definecolor{shade3}{rgb}{1,0.5,0.5}
\definecolor{Emphasis}{rgb}{1,0.5,0.5}
\definecolor{ubuntured}{rgb}{0.788235294,0,0.08627451}
\definecolor{ubuntuorange}{rgb}{1,0.388235294,0.035294118}
\definecolor{ubuntuyellow}{rgb}{1,0.709803922,0.082352941}
\definecolor{paleyellow}{rgb}{1,1,0.8}
\definecolor{black}{rgb}{0,0,0}
\definecolor{RED}{rgb}{1,0,0}
\definecolor{Mulberry}{rgb}{1,1,0.8}

\usepackage{scalefnt}
\usepackage[version=3]{mhchem}

%%%%%%%%%%%%%%%%%%%%%%%%%%%%%%%%%%%%%%%%%%%%%%%%%%%%%%%%%%%%%%%%   Hyphenation

\hyphenation{
AASS
Ab-teil-ung
Ab-teil-ung-en
Alarm-zeichen
aller-dings
Anam-nese
An-eurys-ma
An-eurys-mas
An-eurys-men
An-forder-ung-en
Arbeits-ge-mein-schaft
ARGE
Auf-klär-ung
Aus-bild-ungs-stand-ards
Aus-bildungs-ver-ordnung
Aus-bildungs-zentrum
Aus-wurf-leist-ung
Bauch-aorten-an-eurys-ma
Bauch-aorten-an-eurys-mas
Bauch-aorten-an-eurys-men
Be-at-mung
Be-hand-lungs-mö-glich-keiten
Be-weg-lich-keit
be-wusst-seins-klar
Be-zeich-nung
be-droh-ten
be-reits
Carotis-sono-gra-phie
Des-in-fekt-ion
Dienst-nehmer-haft-pflicht-ge-setz
Donau-stadt
ein-rast-en
ent-weder
Ent-wick-lung
er-ken-nen
Fett-em-bolie
Fing-er-end-glied
Fing-er-end-glied-er
Florids-dorf
Fremd-an-am-nese
Füssig-keits-de-fi-zit
Füssig-keits-de-fi-zits
Füssig-keits-de-fi-zit-es
Ge-lenks-ver-letz-ung
Ge-lenks-ver-letz-ung-en
gleich-ge-stellt
gleich-ge-stellt-en
gleich-zeit-ig
gleich-zeit-ige
gleich-zeit-ig-es
Halb-seiten-läh-mung
Halb-seiten-läh-mung-en
Halb-seiten-schwäche
Haupt-be-schwer-de
Haupt-be-schwer-den
Herz-in-farkt
Herz-in-farkt-es
Herz-in-farkts
hier
hypo-vol-ämi-schen
im-mun-sup-res-sion
im-mun-sup-res-siv
im-mun-sup-res-sive
In-ten-si-tät
Keim-über-träger
Keim-über-träger-in
Keim-über-trägers
Koron-ar-ge-fäße
Körper-kapil-la-ren
Körper-kreis-lauf
Körper-ober-fläche
Kranken-an-stalt
Kranken-an-stalt-en
Kranken-an-stalt-en-ver-bund
Kranken-an-stalt-en-ver-bund-es
Kranken-an-stalt-en-ver-bunds
Kranken-haus
Kranken-haus-es
Kranken-häuser
Kranken-trans-port
Kranken-trans-porte
Kranken-trans-portes
Kranken-trans-ports
Kranken-trans-port-dienst
Kranken-trans-port-dienste
Kranken-trans-port-dienst-es
Kranken-trans-port-wagen
Kranken-trans-port-wag-ens
Kur-an-stalt-en-ge-setz
leicht
Lungen-er-krank-ung
Lungen-er-krank-ung-en
Lungen-kapil-la-re
Lungen-kapil-la-ren
Lungen-kreis-lauf
man-geln-des
Mega-code
Mega-code-train-ing
Mega-code-train-ings
Mehr-höhl-en-ver-letz-ung
Mehr-höhl-en-ver-letz-ung-en
nach-ge-wies-en
nach-ge-wies-en-en
nach-ge-wies-en-es
neuro-lo-gisch
neuro-lo-gische
neuro-lo-gisch-en
neuro-lo-gisch-es
Nie-der-halte-sieb
Nie-der-halte-siebe
Nie-der-halte-siebes
Nie-der-halte-siebs
Not-arzt
Not-arzt-ein-satz-fahr-zeug
Not-arzt-ein-satz-fahr-zeugs
Not-ärzt-in
Not-arzt-hub-schrau-ber
Not-arzt-hub-schrau-bers
Not-arzt-wa-gen
Not-arzt-wa-g-ens
Patient-en-ver-füg-ungs-ge-setz
Per-for-at-ion
Pneu-mo-kok-ken
prä-klin-isch
Puls-oxy-me-ter
Puls-oxy-me-ters
Rea-nim-at-ion
Rea-nim-at-ions-be-reit-schaft
rea-nim-at-ions-pflicht-ig
Rechts-pfle-ge
Rettungs-dienst
Rettungs-trans-port-wagen
Rettungs-trans-port-wag-ens
Sani-tä-ter
Sani-täts-dienst
Schluck-test
schmerz-los-en
Schutz-wirk-ung
sit-uat-ions-ge-recht
sit-uat-ions-ge-rechte
Span-nungs-pneu-mo-thor-ax
Staph-ylo-kok-ken
Symp-tom
symp-tom-los
Tachy-kar-die
Tage
Taschen-atlas
Ter-min-ver-ein-bar-ung
thera-peut-isch
thera-peut-ische
thera-peut-ischen
Tot-raum
un-be-friedi-gend
un-be-friedi-gend-en
Ver-dau-ungs-trakt
Ver-dau-ungs-trakts
Ver-dau-ungs-trakt-es
Ver-nichtungs-schmerz
Vor-hof-flim-mern
Wien
Wien-er
zu-reich-en
}


%%%%%%%%%%%%%%%%%%%%%%%%%%%%%%%%%%%%%%%%%%%%%%%%%
% reset item

 

\usepackage{paralist}
\usepackage{tweaklist}

\usepackage{aass-common}
%%%%%%%%%%%%%%%%%%%%%%%%%%%%%%%%%%%%%%%%%%%%%%%%%%%%%%%%%%%
%%% Symbole
\usepackage{textcomp,marvosym,txfonts,skull}

% resolv symbol clashes 
\usepackage{savesym}
\savesymbol{checkmark}
\savesymbol{Cross} % ifsym
\savesymbol{Sun} % ifsym
\savesymbol{Lightning} % ifsym
\usepackage[geometry,alpine,weather]{ifsym}
\usepackage{dingbat}
\restoresymbol{DINGBAT}{checkmark}



\title[AASS - GER]{Arbeits- und Ausbildungsstandards für den Sanitätsdienst}
\subtitle[Errors]{Gerätelehre}
\author[]{ARGE AASS}
\institute[ARGE AASS]{
Arbeitsgemeinschaft\\
  Arbeits- und Ausbildungsstandards für den Sanitätsdienst\\
  ~\\[1ex]
  \url{http://www.aass.at}
}
\date[November 2004]{November 26, 2004}



\begin{document}

\AtBeginSection[]
{
 \begin{frame}
  \frametitle{Outline}
  \small
  \tableofcontents[currentsection,hideothersubsections]
  \normalsize
 \end{frame}
}

\maketitle


%%%%%%%%%%%%%%%%%%%%%%%%%%%%%%%%%%%%%%%%%%%%%%%%%%%%%%%%%%%
%%%%%%%%%%%%%%%%%%%%%%%%%%%%%%%%%%%%%%%%%%%%%%%%%%%%%%%%%%%
%%%%%%%%%%%%%%%%%%%%%%%%%%%%%%%%%%%%%%%%%%%%%%%%%%%%%%%%%%%
\section{Medizinproduktegesetz - MPG}

%%%%%%%%%%%%%%%%%%%%%%%%%%%%%%%%%%%%%%%%%%%%%%%%%%%%%%%%%%%
\begin{frame}
\frametitle{Das CE-Zeichen}
\centering \includegraphics[width=0.5\linewidth]{./images/geraete/Ce-logo.jpg}
% \gCaptionof{figure}{Das CE-Kennzeichen.}
\end{frame}

\begin{frame}
Here is another slide.
\end{frame}








%%%%%%%%%%%%%%%%%%%%%%%%%%%%%%%%%%%%%%%%%%%%%%%%%%%%%%%%%%%
\begin{frame}
\frametitle{Risikoklassen}
{\tiny
\begin{tabulary}{\linewidth}{lLL}\addlinespace
\gTabSub{Klasse I}
&
Kein oder geringes Anwendungsrisiko
&
z.B: Verbandsmaterial, Gehhilfe, Rollstuhl, Tragsessel, Sehhilfe, Stethoskop, Trage, Beatmungsbeutel, Fieberthermometer, Druckminderer
\\\addlinespace
\gTabSub{Klasse IIa}
&
Anwendungsrisiko ohne entsprechender Ausbildung bzw. Unterweisung
&
z.B: aktive Diagnosegeräte, Absaugeinheit und -katheter, EKG-Kabel, Zuckerstreifen, Trachealtuben, Dauerkatheter
\\\addlinespace
\gTabSub{Klasse IIb}
&
Erhöhtes Anwendungsrisiko ohne entsprechender Ausbildung bzw. Unterweisung
&
z.B: Beatmungsgerät, Defibrillator, Pulsoxymeter, Babyinkubator, Röntgengeräte
\\\addlinespace
\gTabSub{Klasse III}
&
Besonders hohes Anwendungsrisiko ohne entsprechende Ausbildung bzw. Unterweisung
&
z.B: Perfusor, Insulinpumpe, Herzkatheter oder -klappen, Medizinprodukte mit Arzneimittelkomponenten
\\\addlinespace\bottomrule
\end{tabulary}
}

\end{frame}




%%%%%%%%%%%%%%%%%%%%%%%%%%%%%%%%%%%%%%%%%%%%%%%%%%%%%%%%%%%
\begin{frame}
\frametitle{Pflichten des Sanitäters gemäß MPG}

Sanitäter sind lt. \gAbk{MPG} verpflichtet, die
Medizinprodukte richtig und sicher einzusetzen bzw. anzuwenden, so wie es
ihnen bei ihrer Ausbildung bzw. Unterweisung gelehrt wurde.
Außerdem sind sie verpflichtet die Geräte vor dem Gebrauch
mittels eines Gerätecheck auf ihre Funktion zu überprüfen.

\end{frame}
%%%%%%%%%%%%%%%%%%%%%%%%%%%%%%%%%%%%%%%%%%%%%%%%%%%%%%%%%%%
%%%%%%%%%%%%%%%%%%%%%%%%%%%%%%%%%%%%%%%%%%%%%%%%%%%%%%%%%%%
%%%%%%%%%%%%%%%%%%%%%%%%%%%%%%%%%%%%%%%%%%%%%%%%%%%%%%%%%%%
%%%%%%%%%%%%%%%%%%%%%%%%%%%%%%%%%%%%%%%%%%%%%%%%%%%%%%%%%%%
%%%%%%%%%%%%%%%%%%%%%%%%%%%%%%%%%%%%%%%%%%%%%%%%%%%%%%%%%%%
\part{Techniken}
\partpage
 
%%%%%%%%%%%%%%%%%%%%%%%%%%%%%%%%%%%%%%%%%%%%%%%%%%%%%%%%%%%
%%%%%%%%%%%%%%%%%%%%%%%%%%%%%%%%%%%%%%%%%%%%%%%%%%%%%%%%%%%
%%%%%%%%%%%%%%%%%%%%%%%%%%%%%%%%%%%%%%%%%%%%%%%%%%%%%%%%%%%
\section{Blutdruck- und Herzfrequenzmessung}

%%%%%%%%%%%%%%%%%%%%%%%%%%%%%%%%%%%%%%%%%%%%%%%%%%%%%%%%%%%
\begin{frame}
\frametitle{Messung der Herzfrequenz -- HF}

\end{frame}

%%%%%%%%%%%%%%%%%%%%%%%%%%%%%%%%%%%%%%%%%%%%%%%%%%%%%%%%%%%
\begin{frame}
\frametitle{Messung des Blutdrucks nach Riva-Rocci}
\scriptsize
Die Manschette sollte faltenfrei und nicht über Kleidung angelegt werden. Die
Mitte der Luftkammer soll über der Oberarmarterie zu liegen kommen. Manche
Manschetten haben einen Markierungspfeil für die Oberarmarterie, der Pfeil soll
dann oberhalb dieser zu liegen kommen. Man achte auf die Lage der Innen- bzw.
Außenseite!

Nun wird
der Puls an der Radialarterie (A. radialis) getastet. Die Manschette wird
solange aufgepumpt, bis der Puls nicht mehr getastet werden kann. Dann werden
noch 20-30\,mmHg weiter aufgepumpt. Das Stethoskop wird in der Ellenbeuge auf
die Oberarmarterie (A. brachialis, Verlauf eher zur Körpermitte) angelegt. Die
Luft wird langsam aus der Manschette abgelassen. 

Der Wert bei dem man
erstmalig ein Pochen \gZ{(Korotkowsche Geräusche)} hört gilt als
\gEs{systolischer Wert}. Die Luft wird weiter abgelassen. Der Wert, bei dem das
Pochen verschwindet, ist der \gEs{diastolische Wert}. \footnote{Es kann
vorkommen dass das Pochen nie verschwindet sondern nur plötzlich leiser wird.
Dann gilt der Wert bei dem das Pochen \emph{deutlich} leiser wird als
diastolischer Wert.}

Die
Werte werden auf die nächste Fünferstelle abgerundet (z.B.:
122\,/\,82\,mmHg wird abgerundet auf 120\,/\,80\,mmHg). Dokumentiert
wird der Blutdruck mit RR (z.B.: RR\,=\,125\,/\,75\,mmHg). Der
Normalwert beim Erwachsenen liegt ca. bei RR 120\,/\,80\,mmHg
(\gSee{topic:blutdruck}).

\end{frame}

%%%%%%%%%%%%%%%%%%%%%%%%%%%%%%%%%%%%%%%%%%%%%%%%%%%%%%%%%%%
\begin{frame}
\frametitle{Messung des Blutdrucks -- palpatorisch}
\tiny
Die Manschette sollte faltenfrei und nicht über Kleidung angelegt werden. Die
Mitte der Luftkammer soll über der Oberarmarterie zu liegen kommen. Manche
Manschetten haben einen Markierungspfeil für die Oberarmarterie, der Pfeil soll
dann oberhalb dieser zu liegen kommen. Man achte auf die Lage der Innen- bzw.
Außenseite!

Nun wird
der Puls an der Radialarterie (A. radialis) getastet. Die Manschette wird
solange aufgepumpt, bis der Puls nicht mehr getastet werden kann. Dann werden
noch 20-30\,mmHg weiter aufgepumpt. Das Stethoskop wird in der Ellenbeuge auf
die Oberarmarterie (A. brachialis, Verlauf eher zur Körpermitte) angelegt. Die
Luft wird langsam aus der Manschette abgelassen. 

Der Wert bei dem man
erstmalig ein Pochen \gZ{(Korotkowsche Geräusche)} hört gilt als
\gEs{systolischer Wert}. Die Luft wird weiter abgelassen. Der Wert, bei dem das
Pochen verschwindet, ist der \gEs{diastolische Wert}. \footnote{Es kann
vorkommen dass das Pochen nie verschwindet sondern nur plötzlich leiser wird.
Dann gilt der Wert bei dem das Pochen \emph{deutlich} leiser wird als
diastolischer Wert.}

Die
Werte werden auf die nächste Fünferstelle abgerundet (z.B.:
122\,/\,82\,mmHg wird abgerundet auf 120\,/\,80\,mmHg). Dokumentiert
wird der Blutdruck mit RR (z.B.: RR\,=\,125\,/\,75\,mmHg). Der
Normalwert beim Erwachsenen liegt ca. bei RR 120\,/\,80\,mmHg
(\gSee{topic:blutdruck}).
\end{frame}

%%%%%%%%%%%%%%%%%%%%%%%%%%%%%%%%%%%%%%%%%%%%%%%%%%%%%%%%%%%
\begin{frame}
\frametitle{Was ist besonders zu beachten?}
Die Blutdruckmanschette soll nach Möglichkeit nicht auf der Seite
\begin{itemize}
	\item eines Shuntarmes bei Dialysepatienten
	\begin{itemize}
		\item Ein Dialyse-Shunt ist ein künstlicher Kurzschluss
		zwischen einer Armarterie und einer Armvene. Dadurch wird
		die Vene erweitert und sie pulsiert. Mann kann viel mehr
		Durchfluss für die Dialyse erreichen. Die Nahtstelle
		zwischen Vene und Arterie ist jedoch sehr sensibel, wenn
		gestaut wird kann sie aufplatzen $\rightarrow$ Gefahr
		einer massiven Shuntblutung.
	\end{itemize}
		\item eines Armödems, z.B. nach einer
		Brustamputation
		\item einer Halbseitenschwäche oder -lähmung 
\end{itemize}
angelegt werden.

\gCave{Shunt-Arme sind tabu für die Blutdruckmessung!}


\end{frame}

%%%%%%%%%%%%%%%%%%%%%%%%%%%%%%%%%%%%%%%%%%%%%%%%%%%%%%%%%%%
%%%%%%%%%%%%%%%%%%%%%%%%%%%%%%%%%%%%%%%%%%%%%%%%%%%%%%%%%%%
%%%%%%%%%%%%%%%%%%%%%%%%%%%%%%%%%%%%%%%%%%%%%%%%%%%%%%%%%%%
\section{Sauerstoff -- \ce{O2}}
%%%%%%%%%%%%%%%%%%%%%%%%%%%%%%%%%%%%%%%%%%%%%%%%%%%%%%%%%%%
\begin{frame}
\frametitle{Definition}
\gT{Sauerstoff} ist ein
lebenswichtiges Gas welches zu 21\% in unserer Atemluft
vorkommt und außerdem ein essentielles
Notfall{\textquotedblright}medikament{\textquotedblleft}.
Das chemische Symbol lautet \gT{\ce{O2}}.
\end{frame}



\subsection{Sauerstoff -- technische Grundlagen}
%%%%%%%%%%%%%%%%%%%%%%%%%%%%%%%%%%%%%%%%%%%%%%%%%%%%%%%%%%%
\begin{frame}
\frametitle{technische Grundlagen}
Sauerstoff wird in Druckflaschen  gelagert. Laut EU-Norm sind diese
Flaschen weiß und mit einem
{\quotedblbase}\gEs{N}{\textquotedblleft} gekennzeichnet.
Das {\quotedblbase}N{\textquotedblleft} steht für {\quotedblbase}neue Kennzeichnung{\textquotedblleft}, da früher die
Flaschen blau gekennzeichnet wurden. \footnote{{\textquotedblright}N{\textquotedblleft} bedeutet in diesem Fall nicht Stickstoff!}

Die Druckflaschen haben unterschiedliche Füllgrößen (Volumina),
gängige Flaschengrößen sind 0,5, 1, 1,5, 2, 5 und
10\,l. Der Sauerstoff wird unter Druck gelagert, wodurch die
Menge des zur Verfügung stehenden Sauerstoffs um ein Vielfaches
gesteigert wird. Der Druck einer gefüllten Flasche beträgt
normalerweise ca. 200\,bar. 

Jede Druckflasche verfügt über ein \gT{Hauptventil} und
ein Standardgewinde. An diesem Gewinde ist entweder ein
weiterführender Druckschlauch oder der \gE{Druckminderer}
einer \gE{Berieselungseinheit} angeschlossen. 
\end{frame}

%%%%%%%%%%%%%%%%%%%%%%%%%%%%%%%%%%%%%%%%%%%%%%%%%%%%%%%%%%%
\begin{frame}
\frametitle{\gT{Berieselungseinheit}}
Die
Berieselungseinheit besteht im Wesentlichen aus einem
\gT{Druckminderer}. Der Druckminderer hat ein
\gE{Abgabeventil}, welches den \gE{Druck drosselt} und eine Abgabe an den Patienten ermöglicht. Je nach Bauart kann der Druckminderer über ein \gE{Regelventil}, ein \gE{Flowmeter}\index{Flowmeter}
und ein \gE{Manometer}\index{Manometer} verfügen. Mit dem
\gT{Regelventil} kann man die Durchfluss- bzw. Abgaberate ein\-stellen, das
\gT{Flowmeter} (Durchflussanzeige) zeigt die aktuelle
Durchflussrate in Liter pro Minute (\gLMin) an. Das
\gT{Manometer} zeigt den Druck in der Sauerstoffflasche an
wenn das Haupt\-ventil geöffnet ist.
\end{frame}



\subsection{Heiteres Rechnen mit Sauerstoff}
%%%%%%%%%%%%%%%%%%%%%%%%%%%%%%%%%%%%%%%%%%%%%%%%%%%%%%%%%%%
\begin{frame}
\frametitle{Mit Sauerstoff lässt sich gut rechnen ...}



Die insgesamt verfügbare Menge wird errechnet indem man
die Füllgröße der Flasche (in~l) mit dem
Flaschendruck (in bar) multipliziert.

\begin{equation}
\mbox{Verfügbares Volumen in l} = {\begin{array}{c}
\mbox{Volumen der Flasche in l}\\
\mbox{(Füllgröße)} \end{array}}
\times 
\begin{array}{c}{\mbox{Druck in bar}}\\
\mbox{(Flaschendruck)}
\end{array}
\end{equation}

Wenn man wissen möchte, wie lange die vorhandene Sauerstoffmenge bei
einer bestimmten Abgaberate reicht, dividiert man einfach
die gesamt verfügbare Menge durch die Abgaberate. 
\end{frame}

%%%%%%%%%%%%%%%%%%%%%%%%%%%%%%%%%%%%%%%%%%%%%%%%%%%%%%%%%%%
\begin{frame}
\frametitle{}
\gCa{Achtung!} Es ist eine \gEs{ausreichende Reserve}
einzuplanen, in der Flasche sollten min. \gEs{10 bar
Restdruck} verbleiben! 

\gFazit{Es ist eine \gEs{Zeitreserve} und der \gEs{Restdruck} abzuziehen!}
\begin{equation}
\mbox{Zeit} = \frac{\mbox{Verfügbare
Menge}}{\mbox{Abgaberate}}
\end{equation}
\end{frame}

%%%%%%%%%%%%%%%%%%%%%%%%%%%%%%%%%%%%%%%%%%%%%%%%%%%%%%%%%%%
\begin{frame}
\frametitle{Beispiel}
Man hat eine Flasche mit einer Füllgröße von 10\,l. Diese
steht unter einem Druck von 150\,bar: Daraus ergibt
sich, dass in der Flasche 1500\,l \ce{O2} sind.

Der Patient benötigt 5\,\unitfrac{l}{min} \ce{O2} : 

Der Inhalt würde theoretisch für 300 Minuten reichen. Ein
Reservevolumen ist dabei aber noch nicht berücksichtigt!

\end{frame}


\subsection{Umgang mit Sauerstoff und Druckflaschen}
%%%%%%%%%%%%%%%%%%%%%%%%%%%%%%%%%%%%%%%%%%%%%%%%%%%%%%%%%%%
\begin{frame}
\frametitle{Die Sauerstoffflaschen sind pfleglich zu behandeln!}


Sauerstoff ist \gEs{brandfördernd} und zusammen mit Fett
sogar \gEs{selbstentzündlich}! Neben der Brand- und
Explosionsgefahr birgt der große Druck unter dem die Flaschen
stehen weitere Gefahren: Sollte ein Ventil abbrechen kann die
Flasche eine enorme Kraft entwickeln und zu einer Rakete
werden. Die Kraft einer mit 200\,bar gefüllten Flasche reicht
aus um Mauern zu durchbrechen!
\end{frame}

%%%%%%%%%%%%%%%%%%%%%%%%%%%%%%%%%%%%%%%%%%%%%%%%%%%%%%%%%%%
\begin{frame}
\frametitle{Cave}
\gCave{
\gCa{Achtung! Vorsicht beim Hantieren -- Zerknall- und
Explosionsgefahr!}
\begin{enumerate}
    \item Sauerstoff ist zusammen mit Fett selbstentzündlich!
    \begin{enumerate}
        \item Die Armaturen und Gewinde sind fettfrei zu
        halten!
        \item Ventile dürfen nicht mit fetten oder öligen Händen
        bedient werden! 
    \end{enumerate}
    \item Das Hantieren mit Feuer und offenem
    Licht ist in der Nähe von Sauerstoffflaschen verboten!
    \item Die Druckflaschen sind immer zu sichern und vor
    Umfallen, etc. zu schützen!
\end{enumerate}
}
\end{frame}


\subsection{Verabreichungsarten}
%%%%%%%%%%%%%%%%%%%%%%%%%%%%%%%%%%%%%%%%%%%%%%%%%%%%%%%%%%%
\begin{frame}
\frametitle{}
\begin{tabulary}{\linewidth}[H]{CLLL}
\gTabHeading{Was}
 &
\gTabHeading{Wann}
 &
\gTabHeading{Vorteile}
&
\gTabHeading{Nachteile}
\\\addlinespace\toprule\addlinespace 
\gTabSub{Sauerstoff\-brille}
 &
Spontanatmung

\ce{O2}-Flow {\textless}~5\gLMin
 &
gute Toleranz

fehlendes Engegefühl

Sprechen \& Husten möglich
 &
ungenaue Dosierung

Austrocknung  der Schleimhäute 
\\\addlinespace\midrule\addlinespace
\gTabSub{Sauerstoff\-maske}
 &
Spontanatmung

\ce{O2}-Flow ab 5\gLMin
 &
höhere \ce{O2}-Konzentration
 &
\ce{CO2}-Rückatmung  wenn \ce{O2}-Durchfluß
{\textless}~5\gLMin!) 

für Patienten evt. unangenehm bzw. gewöhnungsbedürftig
\\\addlinespace\midrule\addlinespace
\gTabSub{Sauerstoff\-maske mit Reservoir}
 &
Spontanatmung

\ce{O2}-Flow ab 5\gLMin
 &
Noch höhere \ce{O2}-Konzentration
 &
Wie oben.

Reservoir muss zuerst händisch befüllt werden.
\\\addlinespace\bottomrule
\end{tabulary}
\end{frame}

% %%%%%%%%%%%%%%%%%%%%%%%%%%%%%%%%%%%%%%%%%%%%%%%%%%%%%%%%%%%
% \begin{frame}
% \frametitle{Bilderserie: Mittel zur Verabreichung von Sauerstoff}
% \gDias{Bilderserie: Mittel zur Verabreichung von Sauerstoff.}
% {\includegraphics[width=2.5cm]{./images/o2-brille.png}
% \gCaptionof{figure}{Sauerstoffbrille}}
% {\includegraphics[width=2.5cm]{./images/o2-maske.png}
% \gCaptionof{figure}{Sauerstoffmaske}}
% {\includegraphics[width=2.5cm]{./images/o2-maske-reservoir.png}
% \gCaptionof{figure}{Sauerstoffmaske mit Reservoir}}
% \end{frame}


\section{Beatmung}
%%%%%%%%%%%%%%%%%%%%%%%%%%%%%%%%%%%%%%%%%%%%%%%%%%%%%%%%%%%
\begin{frame}
\frametitle{Beatmung}
Bei \gEs{fehlender} oder unzureichender \gEs{Spontanatmung} (Eigenatmung) wird
eine künstliche Beatmung mittels
\gEs{Beatmungsbeutel} (\gSee{topic:beatmungsbeutel}) oder \gEs{Beatmungsgerät}
durchgeführt. Dabei handelt es sich um eine \gE{Überdruckbeatmung}, d.h. es
wird mit Druck Luft in die Atemwege des Patienten gepumpt. (Im Unterscheid dazu
wird bei der natürlichen Atmung Luft mit Unterdruck in den Patienten
eingesogen.) 

Der notwendige Überdruck kann zu Problemen führen, wie z.B. Überschreiten des
Speiseröhren-Öffnungsdrucks (Magenbeatmung, Magenblähung, Erbrechen) oder
Lungenschäden.
\end{frame}

%%%%%%%%%%%%%%%%%%%%%%%%%%%%%%%%%%%%%%%%%%%%%%%%%%%%%%%%%%%
\begin{frame}
\frametitle{Material}
Eine Beatmung durch Fachpersonal sollte wenn möglich immer mit Hilfmitteln wie
dem \gE{Beatmungsbeutel} oder einem \gE{Beatmungsgerät} durchgeführt werden.
Der Beatmungsbeutel ist das einfachste Hilfmittel welches dem Fachpersonal zur
Verfügung steht und kommt in vielen Bereichen (Rettungsdienst, Krankenhaus,
\ldots) zum Einsatz. Er ermöglichst die kontrollierte oder unterstützende
Beatmung eines Patienten.

Wenn vorhanden soll der Beutel mit einer Sauerstoffzufuhr verbunden werden. Am
einfachsten erfolgt dies über einen Verbindungsschlauch und ein eine
\ce{O2}-Berieselungseinheit, welches mit einem Flow von 15\gLMin{}
eingestellt ist. Mit Hilfe eines Reservoirs, welches Sauerstoff während der
Beatmungspausen zwischenspeichert,  kann die \ce{O2}-Konzentration der
Einatemluft weiter gesteigert werden. 

\gZ{Andere \ce{O2}-Systeme sind am Markt
verfügbar (Demand-Ventile, etc.), aber nur begrenzt verbreitet.}
\end{frame}

%%%%%%%%%%%%%%%%%%%%%%%%%%%%%%%%%%%%%%%%%%%%%%%%%%%%%%%%%%%
\begin{frame}
\frametitle{Beatmungsbeutel mit Reservoir, Bakterienfilter und Masken unterschiedlicher Größe.}
\begin{columns}
\begin{column}[T]{0.5\linewidth}
\includegraphics[width=\linewidth]{./images/geraete/ger-ambu.jpg}
\end{column}
\begin{column}[T]{0.5\linewidth}
Beatmungsbeutel mit Reservoir, Bakterienfilter und Masken unterschiedlicher Größe.
\end{column}
\end{columns}
\end{frame}

%%%%%%%%%%%%%%%%%%%%%%%%%%%%%%%%%%%%%%%%%%%%%%%%%%%%%%%%%%%
\begin{frame}
\frametitle{Bestandteile eines
Beatmungsbeutels.}
\begin{columns}
\begin{column}[T]{0.5\linewidth}
\includegraphics[width=\linewidth]{./images/geraete/ambu-komplett-zerlegt.jpg}
\end{column}
\begin{column}[T]{0.5\linewidth}
Ein Beatmungsbeutel besteht im Idealfall aus folgenden Teilen:
\begin{itemize}
	\item Beatmungsmaske 
	\item Bakterienfilter
	\item Ein-/Ausatemventil mit angeschlossenem PEEP-Ventil
	\item Beatmungsbeutel
	\item Reservoir, welches mit einem Verbindungsschlauch
	(\gE{\ce{O2}-Line}) mit einer Berieselungseinheit verbunden
	ist.
\end{itemize}
\end{column}
\end{columns}
\end{frame}


\subsection{Beatmungsgerät}
%%%%%%%%%%%%%%%%%%%%%%%%%%%%%%%%%%%%%%%%%%%%%%%%%%%%%%%%%%%
\begin{frame}
\frametitle{Beatmungsgerät}
\includegraphics[width=7.189cm,height=4.94cm]{./images/beatmungsgeraet-medumat.png}

\captionof{figure}{Beatmungsgerät \emph{Medumat
Standard\textregistered}. \gSourcePicture{ASB
Floridsdorf-Donaustadt}}
\end{frame}

%%%%%%%%%%%%%%%%%%%%%%%%%%%%%%%%%%%%%%%%%%%%%%%%%%%%%%%%%%%
\begin{frame}
\frametitle{}
\ASBFLO{Beim ASB Floridsdorf-Donaustadt kommen folgende Geräte
zum Einsatz:

\begin{itemize}
    \item Hersteller: Fa. Weinmann, Produktfamilie:
    \gT{Medumat} (Notfallbeatmungsgeräte)
    \begin{itemize}
        \item Medumat{\texttrademark} Compact
        \item Medumat{\texttrademark} Standard
        \item Medumat{\texttrademark} Standard a
    \end{itemize}
    \item Hersteller: Fa. Weinmann, Produktfamilie: \gT{Medumat{\texttrademark}}
   
    ((Intensiv-)Transportbe\-atmungs\-ge\-rät\nocite{ThierbachInterhospitaltransfer1})
    \begin{itemize}
        \item Medumat{\texttrademark} Transport
    \end{itemize}
\end{itemize}}
\end{frame}

%%%%%%%%%%%%%%%%%%%%%%%%%%%%%%%%%%%%%%%%%%%%%%%%%%%%%%%%%%%
\begin{frame}
\frametitle{}

Bedienfelder Berieselungseinheit \gE{Modul Oxygen} (li.)
und Beatmungsgerät \gE{Medumat{\texttrademark} Standard} (re.).

\includegraphics[width=\linewidth]{./images/geraete/ger-medumat-standard-bedienfeld.jpg}

\end{frame}

\subsection{PEEP - Ventil}
%%%%%%%%%%%%%%%%%%%%%%%%%%%%%%%%%%%%%%%%%%%%%%%%%%%%%%%%%%%
\begin{frame}
\frametitle{PEEP - Ventil}
\begin{columns}
\begin{column}[T]{0.5\linewidth}
\includegraphics[width=\linewidth]{./images/geraete/peep-ventil.jpg}
\end{column}
\begin{column}[T]{0.5\linewidth}
\gDefinition{PEEP = Positive EndExspiratory Pressure =
positiver end-exspiratorischer Druck}
\end{column}
\end{columns}
\end{frame}

\section{Absauggeräte}
%%%%%%%%%%%%%%%%%%%%%%%%%%%%%%%%%%%%%%%%%%%%%%%%%%%%%%%%%%%
\begin{frame}
\frametitle{}
\begin{tabulary}{\linewidth}{LLLL}
\gTabHeading{Hersteller}
&
\gTabHeading{Produktfamilie}
&
\gTabHeading{Typ}
&
\gTabHeading{Bemerkungen}
\\\addlinespace\midrule
Weinmann
&
\noindent\gT{Accuvac{\texttrademark}}
&
Rescue
&
rote Deckplatte
\\\addlinespace
&
&
Basic
&
grüne Deckplatte
\\\addlinespace
Laerdal
& 
Laerdal Suction Unit{\texttrademark}
&
&
auslaufend
\\\addlinespace
div. Hersteller
&
Handabsaugpumpe
&
RES-Q-VAC
&
\\\addlinespace
div. Hersteller
&
Orosauger
&
&
\\\addlinespace\bottomrule
\end{tabulary}
\end{frame}


%%%%%%%%%%%%%%%%%%%%%%%%%%%%%%%%%%%%%%%%%%%%%%%%%%%%%%%%%%%
\begin{frame}
\frametitle{}
\begin{tabular}{p{\linewidth - 10cm}p{4cm}p{4cm}}
\gTabHeading{Elektrische Absaugeinheit}

Älteres Model von der Fa. Laerdal (Laerdal Suction Unit{\texttrademark})

Neueres Model von der Fa. Weimann (Accuvac{\texttrademark})
&
\vspace{0.1pt}
\includegraphics[width=\linewidth]{./images/geraete/lsu-alt.jpg}
&
\vspace{0.1pt}
\includegraphics[width=\linewidth]{./images/geraete/accuvac.jpg}
\\\midrule
\gTabHeading{Manuelle Absaugpumpen}

Handabsaugpumpe

Fußabsaugpumpe
&
\vspace{0.1pt}
\includegraphics[width=\linewidth]{./images/geraete/handabsaugung.jpg}
&
\vspace{0.1pt}
\includegraphics[width=\linewidth]{./images/geraete/fussabsaugung1.jpg}
\\\midrule
\gTabHeading{Oro-Sauger}

Schleimabsauger für Neugeborene
&
\vspace{0.1pt}
\includegraphics[width=\linewidth]{./images/geraete/orosauger.jpg}
\\\midrule
\gTabHeading{Absaugkatheter}

zum Einführen in Mund oder Nase
&
\vspace{0.1pt}
 \includegraphics[width=4cm]{./images/geraete/absaugkatheter.jpg} \\\bottomrule
\end{tabular}
\end{frame}


%%%%%%%%%%%%%%%%%%%%%%%%%%%%%%%%%%%%%%%%%%%%%%%%%%%%%%%%%%%
\begin{frame}
\frametitle{Absaugkatheter}
\begin{itemize}
    \item Auf Sicht.
    \item Mund: Mundwinkel -- Ohrläppchen.
    \item Nase: Mundwinkel -- Nasenspitze
\end{itemize}
\gCave{Es darf nur auf Sicht abgesaugt werden!}
\end{frame}

\section{Transporttechniken bzw. -geräte}
%%%%%%%%%%%%%%%%%%%%%%%%%%%%%%%%%%%%%%%%%%%%%%%%%%%%%%%%%%%
\begin{frame}
\frametitle{Tragen mit dem Tragering}
\begin{tabularx}{\linewidth}{XXX}
\includegraphics[width=\linewidth]{./images/motal-aass/tragering1b.jpg}
&
\includegraphics[width=\linewidth]{./images/motal-aass/tragering2b.jpg}
&
\includegraphics[width=\linewidth]{./images/motal-aass/tragering3b.jpg}
\\
Tragering &
Transport mit dem Tragering &
Transport mit dem Tragering 
\\\addlinespace\bottomrule
\end{tabularx}
\end{frame}

%%%%%%%%%%%%%%%%%%%%%%%%%%%%%%%%%%%%%%%%%%%%%%%%%%%%%%%%%%%
\begin{frame}
\frametitle{Umgang mit dem Tragsessel}
\begin{columns}
\begin{column}[T]{0.5\linewidth}
\includegraphics[width=\linewidth]{./images/noauth/GER-tragsessel.jpg}
\end{column}
\begin{column}[T]{0.5\linewidth}
Der Patient ist immer anzuschnallen.Beim Transport über
Stiegen hat der Patient immer talwärts zu blicken.
\end{column}
\end{columns}
\end{frame}

%%%%%%%%%%%%%%%%%%%%%%%%%%%%%%%%%%%%%%%%%%%%%%%%%%%%%%%%%%%
\begin{frame}
\frametitle{Tragetuch}
\begin{tabularx}{\linewidth}{XXX}
\includegraphics[width=\linewidth]{./images/motal-aass/tragetuch1b.jpg}
&
\includegraphics[width=\linewidth]{./images/motal-aass/tragetuch2b.jpg}
&
\includegraphics[width=\linewidth]{./images/motal-aass/tragetuch3b.jpg}
\\
Tuch aufbreiten \ldots &
\ldots die näher zum Patient liegende Hälfte
umschlagen\ldots 
& \ldots die obere Hälfte nochmals
umschlagen\ldots 
\\\addlinespace
\includegraphics[width=\linewidth]{./images/motal-aass/tragetuch4b.jpg}
&
\includegraphics[width=\linewidth]{./images/motal-aass/tragetuch5b.jpg}
&
\includegraphics[width=\linewidth]{./images/motal-aass/tragetuch6b.jpg}
\\
\ldots und das Tuch an den Patienten legen. 
&
Patient leicht anheben und Tuch bis zur Hälfte
unterschieben\ldots 
& \ldots Patient schonend zurück legen.
\\\addlinespace
\includegraphics[width=\linewidth]{./images/motal-aass/tragetuch7b.jpg}
&
\includegraphics[width=\linewidth]{./images/motal-aass/tragetuch8b.jpg}
&
\includegraphics[width=\linewidth]{./images/motal-aass/tragetuch9b.jpg}
\\
Von der anderen Seite ebenfalls leicht anheben und
Tuch ausbreiten. 
& 
Zu dritt das Tuch fest anpacken\ldots 
&
\ldots und auf Kommando gemeinsam anheben
\\\addlinespace\bottomrule
\end{tabularx}
\end{frame}

%%%%%%%%%%%%%%%%%%%%%%%%%%%%%%%%%%%%%%%%%%%%%%%%%%%%%%%%%%%
\begin{frame}
\frametitle{}

\end{frame}








\end{document} 